\Chapter{35}

\Verse{1}
{
    \zA{话说宝钗分明听见黛玉克薄他，因惦记着母亲哥哥，并不回头，一径去了。}
}
{
    \eA{Pao ch'ai had, our story goes, distinctly heard Lin Tai-yü's sneer, but in
        her eagerness to see her mother and brother, she did not so much as turn her
        head round, but continued straight on her way.}
}
{
}

\Verse{2}
{
    \zA{这 里黛玉仍旧立于花阴之下，远远的却向怡红院内望着。}
}
{
    \eA{During this time, Lin Tai-yü halted under the shadow of the trees. Upon
        casting a glance, in the distance towards the I Hung Yüan, she observed Li
        Kung-ts'ai,}
}
{
}

\Verse{3}
{
    \zA{只见李纨、迎春、探春、惜 春并丫鬟人等，都向怡红院内去过之后，一起一起的散尽了；只不见凤姐儿来。}
}
{
    \eA{Ying Ch'un, T'an Ch'un, Hsi Ch'un and various inmates wending their steps in
        a body in the direction of the I Hung court; but after they had gone past,
        and company after company of them had dispersed, she only failed to see lady
        Feng come.}
}
{
}

\Verse{4}
{
    \zA{心 里自己盘算说道：“他怎么不来瞧瞧宝玉呢？}
}
{
    \eA{"How is it," she cogitated within herself, "that she doesn't come to see
        Pao-yü? Even supposing that there was some business to detain her,}
}
{
}

\Verse{5}
{
    \zA{便是有事缠住了，他必定也是要来打 个花胡哨，讨老太太、太太的好儿才是呢。}
}
{
    \eA{she should also have put in an appearance, so as to curry favour with our
        venerable senior and Madame Wang. But if she hasn't shown herself at this
        hour of the day, there must certainly be some cause or other."}
}
{
}

\Verse{6}
{
    \zA{今儿这早晚不来，必有原故。”}
}
{
    \eA{While preoccupied with conjectures, she raised her head. At a second glance,}
}
{
}

\Verse{7}
{
    \zA{一面猜 疑，一面抬头再看时，只见花花簇簇一群人，又向怡红院内来了。}
}
{
    \eA{she discerned a crowd of people, as thick as flowers in a bouquet, pursuing
        their way also into the I Hung court. On looking fixedly, she recognised
        dowager lady Chia, leaning on lady Feng's arm,}
}
{
}

\Verse{8}
{
    \zA{定睛看时，却是 贾母搭着凤姐的手，后头邢夫人、王夫人，跟着周姨娘并丫头媳妇等人，都进院去 了。}
}
{
    \eA{followed by Mesdames Hsing and Wang, Mrs. Chou and servant-girls, married
        women and other domestics. In a body they walked into the court. At the
        sight of them, Tai-yü unwittingly nodded her head,}
}
{
}

\Verse{9}
{
    \zA{黛玉看了，不觉点头，想起有父母的好处来，早又泪珠满面。}
}
{
    \eA{and reflected on the benefit of having a father and mother; and tears
        forthwith again bedewed her face. In a while, she beheld Pao-ch'ai, Mrs.
        Hsüeh and the rest likewise go in.}
}
{
}

\Verse{10}
{
    \zA{少顷，只见薛姨 妈、宝钗等也进去了。}
}
{
    \eA{But at quite an unexpected moment she became aware that Tzu Chüan was
        approaching her from behind.}
}
{
}

\Verse{11}
{
    \zA{忽见紫鹃从背后走来，说道：“姑娘吃药去罢，开水又冷了。”}
}
{
    \eA{"Miss," she said, "you had better go and take your medicine! The hot water
        too has got cold." "What do you, after all,}
}
{
}

\Verse{12}
{
    \zA{黛玉道：“你 到底要怎么样？}
}
{
    \eA{mean by keeping on pressing me so?" inquired Tai-yü. "Whether I have it or
        not,}
}
{
}

\Verse{13}
{
    \zA{只是催。}
}
{
    \eA{what's that to you?" "Your cough,"}
}
{
}

\Verse{14}
{
    \zA{我吃不吃，与你什么相干？”}
}
{
    \eA{smiled Tzu Chüan, "has recently got a trifle better, and won't you again
        take your medicine?}
}
{
}

\Verse{15}
{
    \zA{紫鹃笑道：“咳嗽的才好了 些，又不吃药了？}
}
{
    \eA{This is, it's true, the fifth moon, and the weather is hot, but you should,
        nevertheless, take good care of yourself a bit!}
}
{
}

\Verse{16}
{
    \zA{如今虽是五月里，天气热，到底也还该小心些。}
}
{
    \eA{Here you've been at this early hour of the morning standing for ever so long
        in this damp place; so you should go back and have some rest!"}
}
{
}

\Verse{17}
{
    \zA{大清早起，在这 个潮地上站了半日，也该回去歇歇了。”}
}
{
    \eA{This single hint recalled Tai-yü to her senses. She at length realised that
        her legs felt rather tired. After lingering about abstractedly for a long
        while,}
}
{
}

\Verse{18}
{
    \zA{一句话提醒了黛玉，方觉得有点儿腿酸， 呆了半日，方慢慢的扶着紫鹃，回到潇湘馆来。}
}
{
    \eA{she quietly returned into the Hsiao Hsiang lodge, supporting herself on Tzu
        Chüan. As soon as they stepped inside the entrance of the court, her gaze
        was attracted by the confused shadows of the bamboos, which covered the
        ground,}
}
{
}

\Verse{19}
{
    \zA{一进院门，只见满地下竹影参差， 苔痕浓淡，不觉又想起《西厢记》中所云“幽僻处可有人行？}
}
{
    \eA{and the traces of moss, here thick, there thin, and she could not help
        recalling to mind those two lines of the passage in the Hsi Hsiang Chi: "In
        that lone nook some one saunters about, White dew coldly bespecks the
        verdant moss."}
}
{
}

\Verse{20}
{
    \zA{点苍苔白露泠泠”二 句来，因暗暗的叹道：“双文虽然命薄，尚有孀母弱弟；今日我黛玉之薄命，一并 连孀母弱弟俱无。”}
}
{
    \eA{"Shuang Wen," she consequently secretly communed within herself, as she
        sighed, "had of course a poor fate; but she nevertheless had a widowed
        mother and a young brother; but in the unhappy destiny, to which I, Tai-yü,
        am at present doomed, I have neither a widowed mother nor a young brother."}
}
{
}

\Verse{21}
{
    \zA{想到这里，又欲滴下泪来。}
}
{
    \eA{At this point in her reflections, she was about to melt into another fit of
        crying,}
}
{
}

\Verse{22}
{
    \zA{不防廊下的鹦哥见黛玉来了，“嘎” 的一声扑了下来，倒吓了一跳。}
}
{
    \eA{when of a sudden, the parrot under the verandah caught sight of Tai-yü
        approaching, and, with a shriek, he jumped down from his perch, and made her
        start with fright. "Are you bent upon compassing your own death!"}
}
{
}

\Verse{23}
{
    \zA{因说道：“你作死呢，又？}
}
{
    \eA{she exclaimed. "You've covered my head all over with dust again!"}
}
{
}

\Verse{24}
{
    \zA{了我一头灰。”}
}
{
    \eA{The parrot flew back to his perch.}
}
{
}

\Verse{25}
{
    \zA{那鹦哥 又飞上架去，便叫：“雪雁，快掀帘子，姑娘来了！”}
}
{
    \eA{"Hsüeh Yen," he kept on shouting, "quick, raise the portiere! Miss is come!"
        Tai-yü stopped short and rapped on the frame with her hand.}
}
{
}

\Verse{26}
{
    \zA{黛玉便止住步，以手扣架， 道：“添了食水不曾？”}
}
{
    \eA{"Have his food and water been replenished?" she asked. The parrot forthwith
        heaved a deep sigh, closely resembling, in sound,}
}
{
}

\Verse{27}
{
    \zA{那鹦哥便长叹一声，竟大似黛玉素日吁嗟音韵，接着念道： “侬今葬花人笑痴，他年葬侬知是谁！”}
}
{
    \eA{the groans usually indulged in by Tai-yü, and then went on to recite: "Here
        I am fain these flowers to inter, but humankind will laugh me as a fool."
        Who knows who will in years to come commit me to my grave. As soon as these
        lines fell on the ear of Tai-yü and Tzu Chüan,}
}
{
}

\Verse{28}
{
    \zA{黛玉紫鹃听了，都笑起来。}
}
{
    \eA{they blurted out laughing. "This is what you were repeating some time back,}
}
{
}

\Verse{29}
{
    \zA{紫鹃笑道：“这 都是素日姑娘念的，难为他怎么记了。”}
}
{
    \eA{Miss." Tzu Chüan laughed, "How did he ever manage to commit it to memory?"
        Tai-yü then directed some one to take down the frame and suspend it instead
        on a hook,}
}
{
}

\Verse{30}
{
    \zA{黛玉便命将架摘下来另挂在月洞窗外的钩 上。}
}
{
    \eA{outside the circular window, and presently entering her room, she seated
        herself inside the circular window.}
}
{
}

\Verse{31}
{
    \zA{于是进了屋子，在月洞窗内坐了，吃毕药。}
}
{
    \eA{She had just done drinking her medicine, when she perceived that the shade
        cast by the cluster of bamboos,}
}
{
}

\Verse{32}
{
    \zA{只见窗外竹影映入纱窗，满屋内阴 阴翠润，几簟生凉。}
}
{
    \eA{planted outside the window, was reflected so far on the gauze lattice as to
        fill the room with a faint light, so green and mellow, and to impart a
        certain coolness to the teapoys and mats.}
}
{
}

\Verse{33}
{
    \zA{黛玉无可释闷，便隔着纱窗，调逗鹦哥做戏，又将素日所喜的 诗词也教与他念。}
}
{
    \eA{But Tai-yü had no means at hand to dispel her ennui, so from inside the
        gauze lattice, she instigated the parrot to perform his pranks; and
        selecting some verses, which had ever found favour with her, she tried to
        teach them to him.}
}
{
}

\Verse{34}
{
    \zA{这且不在话下。}
}
{
    \eA{But without descending to particulars,}
}
{
}

\Verse{35}
{
    \zA{且说宝钗来至家中，只见母亲正梳头呢，看见他进来，便笑着说道：“你这么 早就梳上头了。”}
}
{
    \eA{let us now advert to Hsüeh Pao-ch'ai. On her return home, she found her
        mother alone combing her hair and having a wash. "Why do you run over at
        this early hour of the morning?" she speedily inquired when she saw her
        enter.}
}
{
}

\Verse{36}
{
    \zA{宝钗道：“我瞧瞧妈妈身上好不好。}
}
{
    \eA{"To see," replied Pao-ch'ai, "whether you were all right or not, mother. Did
        he come again,}
}
{
}

\Verse{37}
{
    \zA{昨儿我去了，不知他可又过 来闹了没有？”}
}
{
    \eA{I wonder, after I left yesterday and make any more trouble or not?" As she
        spoke, she sat by her mother's side, but unable to curb her tears,}
}
{
}

\Verse{38}
{
    \zA{一面说，一面在他母亲身旁坐下，由不得哭将起来。}
}
{
    \eA{she began to weep. Seeing her sobbing, Mrs. Hsüeh herself could not check
        her feelings, and she, too, burst out into a fit of crying.}
}
{
}

\Verse{39}
{
    \zA{薛姨妈见他一 哭，自己掌不住也就哭了一场，一面又劝他：“我的儿，你别委屈了。}
}
{
    \eA{"My child," she simultaneously exhorted her, "don't feel aggrieved! Wait,
        and I'll call that child of wrath to order; for were anything to happen to
        you, from whom will I have anything to hope?"}
}
{
}

\Verse{40}
{
    \zA{你等我处分 那孽障。}
}
{
    \eA{Hsüeh P'an was outside and happened to overhear their conversation,}
}
{
}

\Verse{41}
{
    \zA{你要有个好歹，叫我指望那一个呢？”}
}
{
    \eA{so with alacrity he ran over, and facing Pao-ch'ai he made a bow, now to the
        left and now to the right,}
}
{
}

\Verse{42}
{
    \zA{薛蟠在外听见，连忙的跑过来，对 着宝钗左一个揖右一个揖，只说：“好妹妹恕我这次罢!原是我昨儿吃了酒，回来
        的晚了，路上撞客着了，来家没醒，不知胡说了些什么，连自己也不知道，怨不得 你生气。”}
}
{
    \eA{observing the while: "My dear sister, forgive me this time. The fact is that
        I took some wine yesterday; I came back late, as I met a few friends on the
        way. On my return home, I hadn't as yet got over the fumes, so I
        unintentionally talked a lot of nonsense. But I don't so much as remember
        anything about all I said.}
}
{
}

\Verse{43}
{
    \zA{宝钗原是掩面而哭，听如此说由不得也笑了，遂抬头向地下啐了一口， 说道：“你不用做这些像生儿了。}
}
{
    \eA{It isn't worth your while, however, losing your temper over such a thing!"
        Pao-ch'ai was, in fact, weeping, as she covered her face, but the moment
        this language fell on her ear, she could scarcely again refrain from
        laughing.}
}
{
}

\Verse{44}
{
    \zA{我知道你的心里多嫌我们娘儿们，你是变着法儿 叫我们离了你就心净了。”}
}
{
    \eA{Forthwith raising her head, she sputtered contemptuously on the ground. "You
        can well dispense with all this sham!" she exclaimed, "I'm well aware that
        you so dislike us both,}
}
{
}

\Verse{45}
{
    \zA{薛蟠听说，连忙笑道：“妹妹这从那里说起？}
}
{
    \eA{that you're anxious to devise some way of inducing us to part company with
        you, so that you may be at liberty."}
}
{
}

\Verse{46}
{
    \zA{妹妹从来不是这么多心说歪话的 人哪。”}
}
{
    \eA{Hsüeh P'an, at these words, hastened to smile. "Sister," he argued, "what
        makes you say so?}
}
{
}

\Verse{47}
{
    \zA{薛姨妈忙又接着道：“你只会听你妹妹的‘歪话’，难道昨儿晚上你说的 那些话，就使得吗？}
}
{
    \eA{once upon a time, you weren't so suspicious and given to uttering anything
        so perverse!" Mrs. Hsüeh hurriedly took up the thread of the conversation.
        "All you know," she interposed, "is to find fault with your sister's remarks
        as being perverse;}
}
{
}

\Verse{48}
{
    \zA{当真是你发昏了？”}
}
{
    \eA{but can it be that what you said last night was the proper thing to say?}
}
{
}

\Verse{49}
{
    \zA{薛蟠道：“妈妈也不必生气，妹妹也不用 烦恼，从今以后，我再不和他们一块儿喝酒了。}
}
{
    \eA{In very truth, you were drunk!" "There's no need for you to get angry,
        mother!" Hsüeh P'an rejoined, "nor for you sister either; for from this day,
        I shan't any more make common cause with them nor drink wine or gad about.}
}
{
}

\Verse{50}
{
    \zA{好不好？”}
}
{
    \eA{What do you say to that?"}
}
{
}

\Verse{51}
{
    \zA{宝钗笑道：“这才明白 过来了。”}
}
{
    \eA{"That's equal to an acknowledgment of your failings," Pao-ch'ai laughed.
        "Could you exercise such strength of will,"}
}
{
}

\Verse{52}
{
    \zA{薛姨妈道：“你要有个横劲，那龙也下蛋了。”}
}
{
    \eA{added Mrs. Hsüeh, "why, the dragon too would lay eggs." "If I again go and
        gad about with them," Hsüeh P'an replied, "and you, sister, come to hear of
        it,}
}
{
}

\Verse{53}
{
    \zA{薛蟠道：“我要再和他 们一处喝，妹妹听见了，只管啐我，再叫我畜生、不是人如何？}
}
{
    \eA{you can freely spit in my face and call me a beast and no human being. Do
        you agree to that? But why should you two be daily worried; and all through
        me alone? For you, mother, to be angry on my account is anyhow excusable;}
}
{
}

\Verse{54}
{
    \zA{何苦来为我一个人， 娘儿两个天天儿操心。}
}
{
    \eA{but for me to keep on worrying you, sister, makes me less then ever worthy
        of the name of a human being!}
}
{
}

\Verse{55}
{
    \zA{妈妈为我生气还犹可，要只管叫妹妹为我操心，我更不是人 了。}
}
{
    \eA{If now that father is no more, I manage, instead of showing you plenty of
        filial piety, mamma, and you, sister, plenty of love, to provoke my mother
        to anger,}
}
{
}

\Verse{56}
{
    \zA{如今父亲没了，我不能多孝顺妈妈，多疼妹妹，反叫娘母子生气、妹妹烦恼， 连个畜生不如了！”}
}
{
    \eA{and annoy my sister, why I can't compare myself to even a four-footed
        creature!" While from his mouth issued these words, tears rolled down from
        his eyes; for he too found it hard to contain them. Mrs. Hsüeh had not at
        first been overcome by her feelings;}
}
{
}

\Verse{57}
{
    \zA{口里说着，眼睛里掌不住掉下泪来。}
}
{
    \eA{but the moment his utterances reached her ear, she once more began to
        experience the anguish,}
}
{
}

\Verse{58}
{
    \zA{薛姨妈本不哭了，听他一 说又伤起心来。}
}
{
    \eA{which they stirred in her heart. Pao-ch'ai made an effort to force a smile.
        "You've already," she said,}
}
{
}

\Verse{59}
{
    \zA{宝钗勉强笑道：“你闹够了，这会子又来招着妈妈哭了。”}
}
{
    \eA{"been the cause of quite enough trouble, and do you now provoke mother to
        have another cry?" Hearing this, Hsüeh P'an promptly checked his tears.}
}
{
}

\Verse{60}
{
    \zA{薛蟠听 说，忙收泪笑道：“我何曾招妈妈哭来着？}
}
{
    \eA{As he put on a smiling expression, "When did I," he asked, "make mother cry?
        But never mind; enough of this!}
}
{
}

\Verse{61}
{
    \zA{罢罢罢，扔下这个别提了，叫香菱来倒 茶妹妹喝。”}
}
{
    \eA{let's drop the matter, and not allude to it any more! Call Hsiang Ling to
        come and give you a cup of tea, sister!" "I don't want any tea."}
}
{
}

\Verse{62}
{
    \zA{宝钗道：“我也不喝茶，等妈妈洗了手，我们就进去了。”}
}
{
    \eA{Pao-ch'ai answered. "I'll wait until mother has finished washing her hands
        and then go with her into the garden." "Let me see your necklet, sister,"
        Hsüeh P'an continued.}
}
{
}

\Verse{63}
{
    \zA{薛蟠道： “妹妹的项圈我瞧瞧，只怕该炸一炸去了。”}
}
{
    \eA{"I think it requires cleaning." "It is so yellow and bright," rejoined
        Pao-ch'ai, "and what's the use of cleaning it again?"}
}
{
}

\Verse{64}
{
    \zA{宝钗道：“黄澄澄的，又炸他做什么？”}
}
{
    \eA{"Sister," proceeded Hsüeh P'an, "you must now add a few more clothes to your
        wardrobe,}
}
{
}

\Verse{65}
{
    \zA{薛蟠又道：“妹妹如今也该添补些衣裳了，要什么颜色花样，告诉我。”}
}
{
    \eA{so tell me what colour and what design you like best." "I haven't yet worn
        out all the clothes I have," Pao-ch'ai explained, "and why should I have
        more made?"}
}
{
}

\Verse{66}
{
    \zA{宝钗道： “连那些衣裳我还没穿遍了，又做什么？”}
}
{
    \eA{But, in a little time, Mrs. Hsüeh effected the change in her costume, and
        hand in hand with Pao-ch'ai, she started on her way to the garden.}
}
{
}

\Verse{67}
{
    \zA{一时薛姨妈换了衣裳，拉着宝钗进去， 薛蟠方出去了。}
}
{
    \eA{Hsüeh P'an thereupon took his departure. During this while, Mrs. Hsüeh and
        Pao-ch'ai trudged in the direction of the garden to look up Pao-yü.}
}
{
}

\Verse{68}
{
    \zA{这里薛姨妈和宝钗进园来看宝玉。}
}
{
    \eA{As soon as they reached the interior of the I Hung court,}
}
{
}

\Verse{69}
{
    \zA{到了怡红院中，只见抱厦里外回廊上许多丫 头老婆站着，便知贾母等都在这里。}
}
{
    \eA{they saw a large concourse of waiting-maids and matrons standing inside as
        well as outside the antechambers and they readily concluded that old lady
        Chia and the other ladies were assembled in his rooms.}
}
{
}

\Verse{70}
{
    \zA{母女两个进来，大家见过了。}
}
{
    \eA{Mrs. Hsüeh and her daughter stepped in. After exchanging salutations with
        every one present,}
}
{
}

\Verse{71}
{
    \zA{只见宝玉躺在榻 上，薛姨妈问他：“可好些？”}
}
{
    \eA{they noticed that Pao-yü was reclining on the couch and Mrs. Hsüeh inquired
        of him whether he felt any better. Pao-yü hastily attempted to bow. "I'm
        considerably better;"}
}
{
}

\Verse{72}
{
    \zA{宝玉忙欲欠身，口里答应着：“好些。”}
}
{
    \eA{he said. "All I do," he went on, "is to disturb you, aunt, and you, my
        cousin, but I don't deserve such attentions."}
}
{
}

\Verse{73}
{
    \zA{又说：“只 管惊动姨娘姐姐，我当不起。”}
}
{
    \eA{Mrs. Hsüeh lost no time in supporting and laying him down. "Mind you tell me
        whatever may take your fancy!"}
}
{
}

\Verse{74}
{
    \zA{薛姨妈忙扶他睡下，又问他：“想什么，只管告诉 我。”}
}
{
    \eA{she proceeded. "If I do fancy anything," retorted Pao-yü smilingly, "I shall
        certainly send to you, aunt, for it." "What would you like to eat,"}
}
{
}

\Verse{75}
{
    \zA{宝玉笑道：“我想起来，自然和姨娘要去。”}
}
{
    \eA{likewise inquired Madame Wang, "so that I may, on my return, send it round
        to you?" "There's nothing that I care for,"}
}
{
}

\Verse{76}
{
    \zA{王夫人又问：“你想什么吃？}
}
{
    \eA{smiled Pao-yü, "though the soup made for me the other day, with young lotus
        leaves,}
}
{
}

\Verse{77}
{
    \zA{回来好给你送来。”}
}
{
    \eA{and small lotus cores was, I thought, somewhat nice."}
}
{
}

\Verse{78}
{
    \zA{宝玉笑道：“也倒不想什么吃。}
}
{
    \eA{"From what I hear, its flavour is nothing very grand," lady Feng chimed in
        laughingly,}
}
{
}

\Verse{79}
{
    \zA{倒是那一回做的那小荷叶儿小 莲蓬儿的汤还好些。”}
}
{
    \eA{from where she stood on one side. "It involves, however, a good deal of
        trouble to concoct; and here you deliberately go and fancy this very thing."}
}
{
}

\Verse{80}
{
    \zA{凤姐一旁笑道：“都听听!口味倒不算高贵，只是太磨牙了。}
}
{
    \eA{"Go and get it ready!" cried dowager lady Chia several successive times.
        "Venerable ancestor," urged lady Feng with a smile, "don't you bother
        yourself about it!}
}
{
}

\Verse{81}
{
    \zA{巴巴儿的想这个吃！”}
}
{
    \eA{Let me try and remember who can have put the moulds away!"}
}
{
}

\Verse{82}
{
    \zA{贾母便一叠连声的叫做去。}
}
{
    \eA{Then turning her head round, "Go and bid," she enjoined an old matron,}
}
{
}

\Verse{83}
{
    \zA{凤姐笑道：“老祖宗别急，我想 想这模子是谁收着呢？”}
}
{
    \eA{"the chief in the cook-house go and apply for them!" After a considerable
        lapse of time, the matron returned. "The chief in the cook-house,"}
}
{
}

\Verse{84}
{
    \zA{因回头吩咐个老婆问管厨房的去要。}
}
{
    \eA{she explained, "says that the four sets of moulds for soups have all been
        handed up."}
}
{
}

\Verse{85}
{
    \zA{那老婆去了半天，来 回话：“管厨房的说：‘四副汤模子都缴上来了。’”}
}
{
    \eA{Upon hearing this, lady Feng thought again for a while. "Yes, I remember,"
        she afterwards remarked, "they were handed up, but I can't recollect to whom
        they were given.}
}
{
}

\Verse{86}
{
    \zA{凤姐听说，又想了一想道： “我也记得交上来了，就只不记得交给谁了。}
}
{
    \eA{Possibly they're in the tea-room." Thereupon, she also despatched a servant
        to go and inquire of the keeper of the tea-room about them; but he too had
        not got them; and it was subsequently the butler,}
}
{
}

\Verse{87}
{
    \zA{多半是在茶房里。”}
}
{
    \eA{entrusted with the care of the gold and silver articles,}
}
{
}

\Verse{88}
{
    \zA{又遣人去问管茶 房的，也不曾收。}
}
{
    \eA{who brought them round. Mrs. Hsüeh was the first to take them and examine
        them. What, in fact, struck her gaze was a small box,}
}
{
}

\Verse{89}
{
    \zA{次后还是管金银器的送了来了。}
}
{
    \eA{the contents of which were four sets of silver moulds. Each of these was
        over a foot long,}
}
{
}

\Verse{90}
{
    \zA{薛姨妈先接过来瞧时，原来是个小匣子，里面装着四副银模子，都有一尺多长， 一寸见方。}
}
{
    \eA{and one square inch (in breadth). On the top, holes were bored of the size
        of beans. Some resembled chrysanthemums, others plum blossom. Some were in
        the shape of lotus seed-cases, others like water chestnuts.}
}
{
}

\Verse{91}
{
    \zA{上面凿着豆子大小，也有菊花的，也有梅花的，也有莲蓬的，也有菱角 的：共有三四十样，打的十分精巧。}
}
{
    \eA{They numbered in all thirty or forty kinds, and were ingeniously executed.
        "In your mansion," she felt impelled to observe smilingly to old lady Chia
        and Madame Wang, "everything has been amply provided for!}
}
{
}

\Verse{92}
{
    \zA{因笑向贾母王夫人道：“你们府上也都想绝了， 吃碗汤还有这些样子。}
}
{
    \eA{Have you got all these things to prepare a plate of soup with! Hadn't you
        told me, and I happened to see them, I wouldn't have been able to make out
        what they were intended for!"}
}
{
}

\Verse{93}
{
    \zA{要不说出来，我见了这个，也不认得是做什么用的。”}
}
{
    \eA{Lady Feng did not allow time to any one to put in her word. "Aunt," she
        said, "how could you ever have divined that these were used last year for
        the imperial viands!}
}
{
}

\Verse{94}
{
    \zA{凤姐 儿也不等人说话，便笑道：“姑妈不知道：这是旧年备膳的时候儿，他们想的法儿。}
}
{
    \eA{They thought of a way by which they devised, somehow or other, I can't tell
        how, some dough shapes, which borrow a little of the pure fragrance of the
        new lotus leaves. But as all mainly depends upon the quality of the soup,}
}
{
}

\Verse{95}
{
    \zA{不知弄什么面印出来，借点新荷叶的清香，全仗着好汤，我吃着究竟也没什么意思。}
}
{
    \eA{they're not, after all, of much use! Yet who often goes in for such soup! It
        was made once only, and that at the time when the moulds were brought; and
        how is it that he has come to think of it to-day?"}
}
{
}

\Verse{96}
{
    \zA{谁家长吃他？}
}
{
    \eA{So speaking, she took (the moulds),}
}
{
}

\Verse{97}
{
    \zA{那一回呈样做了一回，他今儿怎么想起来了！”}
}
{
    \eA{and handed them to a married woman, to go and issue directions to the people
        in the cook-house to procure at once several fowls,}
}
{
}

\Verse{98}
{
    \zA{说着，接过来递与个 妇人，吩咐厨房里立刻拿几只鸡，另外添了东西，做十碗汤来。}
}
{
    \eA{and to add other ingredients besides and prepare ten bowls of soup. "What do
        you want all that lot for?" observed Madame Wang. "There's good reason for
        it," answered lady Feng. "A dish of this kind isn't,}
}
{
}

\Verse{99}
{
    \zA{王夫人道：“要这 些做什么？”}
}
{
    \eA{at ordinary times, very often made, and were, now that brother Pao-yü has
        alluded to it,}
}
{
}

\Verse{100}
{
    \zA{凤姐笑道：“有个原故：这一宗东西家常不大做，今儿宝兄弟提起来 了，单做给他吃，老太太、姑妈、太太都不吃，似乎不大好。}
}
{
    \eA{only sufficient prepared for him, and none for you, dear senior, you, aunt,
        and you, Madame Wang, it won't be quite the thing! So isn't it better that
        this opportunity should be availed of to get ready a whole supply so that
        every one should partake of some, and that even I should, through my
        reliance on your kind favour,}
}
{
}

\Verse{101}
{
    \zA{不如就势儿弄些大家 吃吃，托赖着连我也尝个新儿。”}
}
{
    \eA{taste this novel kind of relish." "You are sharper than a monkey!" Dowager
        lady Chia laughingly exclaimed in reply to her proposal.}
}
{
}

\Verse{102}
{
    \zA{贾母听了，笑道：“猴儿，把你乖的!拿着官中 的钱做人情。”}
}
{
    \eA{"You make use of public money to confer boons upon people." This remark
        evoked general laughter. "This is a mere bagatelle!" eagerly laughed lady
        Feng.}
}
{
}

\Verse{103}
{
    \zA{说的大家笑了。}
}
{
    \eA{"Even I can afford to stand you such a small treat!"}
}
{
}

\Verse{104}
{
    \zA{凤姐忙笑道：“这不相干。}
}
{
    \eA{Then turning her head round, "Tell them in the cook-house,"}
}
{
}

\Verse{105}
{
    \zA{这个小东道儿我还孝敬 的起。”}
}
{
    \eA{she said to a married woman, "to please make an extra supply, and that
        they'll get the money from me."}
}
{
}

\Verse{106}
{
    \zA{便回头吩咐妇人：“说给厨房里，只管好生添补着做了，在我帐上领银子。”}
}
{
    \eA{The matron assented and went out of the room. Pao-ch'ai, who was standing
        near, thereupon interposed with a smile. "During the few years that have
        gone by since I've come here,}
}
{
}

\Verse{107}
{
    \zA{婆子答应着去了。}
}
{
    \eA{I've carefully noticed that sister-in-law Secunda,}
}
{
}

\Verse{108}
{
    \zA{宝钗一旁笑道：“我来了这么几年，留神看起来，二嫂子凭他怎么巧，再巧不 过老太太。”}
}
{
    \eA{cannot, with all her acumen, outwit our venerable ancestor." "My dear
        child!" forthwith replied old lady Chia at these words. "I'm now quite an
        old woman, and how can there still remain any wit in me!}
}
{
}

\Verse{109}
{
    \zA{贾母听说，便答道：“我的儿!我如今老了，那里还巧什么？}
}
{
    \eA{When I was, long ago, of your manlike cousin Feng's age, I had far more wits
        about me than she has! Albeit she now avers that she can't reach our
        standard,}
}
{
}

\Verse{110}
{
    \zA{当日我像 凤丫头这么大年纪，比他还来得呢。}
}
{
    \eA{she's good enough; and compared with your aunt Wang, why, she's infinitely
        superior. Your aunt, poor thing, won't speak much!}
}
{
}

\Verse{111}
{
    \zA{他如今虽说不如我，也就算好了，比你姨娘强 远了!你姨娘可怜见的，不大说话，和木头似的，公婆跟前就不献好儿。}
}
{
    \eA{She's like a block of wood; and when with her father and mother-in-law, she
        won't show herself off to advantage. But that girl Feng has a sharp tongue,
        so is it a wonder if people take to her." "From what you say," insinuated
        Pao-yü with a smile,}
}
{
}

\Verse{112}
{
    \zA{凤儿嘴乖， 怎么怨得人疼他。”}
}
{
    \eA{"those who don't talk much are not loved." "Those who don't speak much,"}
}
{
}

\Verse{113}
{
    \zA{宝玉笑道：“要这么说，不大说话的就不疼了？”}
}
{
    \eA{resumed dowager lady Chia, "possess the endearing quality of reserve. But
        among those, with glib tongues,}
}
{
}

\Verse{114}
{
    \zA{贾母道：“不 大说话的，又有不大说话的可疼之处。}
}
{
    \eA{there's also a certain despicable lot; thus it's better, in a word, not to
        have too much to say for one's self."}
}
{
}

\Verse{115}
{
    \zA{嘴乖的也有一宗可嫌的，倒不如不说的好。”}
}
{
    \eA{"Quite so," smiled Pao-yü, "yet though senior sister-in-law Chia Chu
        doesn't, I must confess, talk much,}
}
{
}

\Verse{116}
{
    \zA{宝玉笑道：“这就是了。}
}
{
    \eA{you, venerable ancestor, treat her just as you do cousin Feng.}
}
{
}

\Verse{117}
{
    \zA{我说大嫂子倒不大说话呢，老太太也是和凤姐姐一样的疼。}
}
{
    \eA{But if you maintain that those alone, who can talk, are worthy of love, then
        among all these young ladies,}
}
{
}

\Verse{118}
{
    \zA{要说单是会说话的可疼，这些姐妹里头也只凤姐姐和林妹妹可疼了。”}
}
{
    \eA{sister Feng and cousin Lin are the only ones good enough to be loved." "With
        regard to the young ladies," remarked dowager lady Chia,}
}
{
}

\Verse{119}
{
    \zA{贾母道：“提 起姐妹，不是我当着姨太太的面奉承：千真万真，从我们家里四个女孩儿算起，都 不如宝丫头。”}
}
{
    \eA{"it isn't that I have any wish to flatter your aunt Hsüeh in her presence,
        but it is a positive and incontestable fact that there isn't, beginning from
        the four girls in our household, a single one able to hold a candle to that
        girl Pao-ch'ai."}
}
{
}

\Verse{120}
{
    \zA{薛姨妈听了，忙笑道：“这话是老太太说偏了。”}
}
{
    \eA{At these words, Mrs. Hsüeh promptly smiled. "Dear venerable senior!" she
        said, "you're rather partial in your verdict."}
}
{
}

\Verse{121}
{
    \zA{王夫人忙又笑道： “老太太时常背地里和我说宝丫头好，这倒不是假说。”}
}
{
    \eA{"Our dear senior," vehemently put in Madame Wang, also smiling, "has often
        told me in private how nice your daughter Pao-ch'ai is; so this is no lie."
        Pao-yü had tried to lead old lady Chia on,}
}
{
}

\Verse{122}
{
    \zA{宝玉勾着贾母，原为要赞 黛玉，不想反赞起宝钗来，倒也意出望外，便看着宝钗一笑。}
}
{
    \eA{originally with the idea of inducing her to speak highly of Lin Tai-yü, but
        when unawares she began to eulogise Pao-ch'ai instead the result exceeded
        all his thoughts and went far beyond his expectations. Forthwith he cast a
        glance at Pao-chai,}
}
{
}

\Verse{123}
{
    \zA{宝钗早扭过头去和袭 人说话去了。}
}
{
    \eA{and gave her a smile, but Pao-chai at once twisted her head round and went
        and chatted with Hsi Jen.}
}
{
}

\Verse{124}
{
    \zA{忽有人来请吃饭，贾母方立起身来，命宝玉：“好生养着罢。”}
}
{
    \eA{But of a sudden, some one came to ask them to go and have their meal.
        Dowager lady Chia rose to her feet, and enjoined Pao-yü to be careful of
        himself.}
}
{
}

\Verse{125}
{
    \zA{把丫头们又嘱 咐了一回，方扶着凤姐儿，让着薛姨妈，大家出房去了。}
}
{
    \eA{She then gave a few directions to the waiting-maids, and resting her weight
        on lady Feng's arm, and pressing Mrs. Hsüeh to go out first, she, and all
        with her, left the apartment in a body.}
}
{
}

\Verse{126}
{
    \zA{犹问：“汤好了不曾？”}
}
{
    \eA{But still she kept on inquiring whether the soup was ready or not.}
}
{
}

\Verse{127}
{
    \zA{又问薛姨妈等：“想什么吃，只管告诉我，我有本事叫凤丫头弄了来咱们吃。”}
}
{
    \eA{"If there's anything you might fancy to eat," she also said to Mrs. Hsüeh
        and the others, "mind you, come and tell me, and I know how to coax that
        hussey Feng to get it for you as well as me."}
}
{
}

\Verse{128}
{
    \zA{薛 姨妈笑道：“老太太也会怄他，时常他弄了东西来孝敬，究竟又吃不多儿。”}
}
{
    \eA{"My venerable senior!" rejoined Mrs. Hsüeh, "you do have the happy knack of
        putting her on her mettle; but though she has often got things ready for
        you, you've, after all, not eaten very much of them."}
}
{
}

\Verse{129}
{
    \zA{凤姐 儿笑道：“姑妈倒别这么说。}
}
{
    \eA{"Aunt," smiled lady Feng, "don't make such statements!}
}
{
}

\Verse{130}
{
    \zA{我们老祖宗只是嫌人肉酸，要不嫌人肉酸，早已把我 还吃了呢！”}
}
{
    \eA{If our worthy senior hasn't eaten me up it's purely and simply because she
        dislikes human flesh as being sour. Did she not look down upon it as sour,
        why, she would long ago have gobbled me up!"}
}
{
}

\Verse{131}
{
    \zA{一句话没说了，引的贾母众人都哈哈的大笑起来。}
}
{
    \eA{This joke was scarcely ended, when it so tickled the fancy of old lady Chia
        and all the inmates that they broke out with one voice in a boisterous fit
        of laughter.}
}
{
}

\Verse{132}
{
    \zA{宝玉在屋里也掌不 住笑了。}
}
{
    \eA{Even Pao-yü, who was inside the room, could not keep quiet. "Really,"}
}
{
}

\Verse{133}
{
    \zA{袭人笑道：“真真的二奶奶的嘴，怕死人。”}
}
{
    \eA{Hsi Jen laughed, "the mouth of our mistress Secunda is enough to terrify
        people to death!"}
}
{
}

\Verse{134}
{
    \zA{宝玉伸手拉着袭人笑道：“你站了这半日，可乏了。”}
}
{
    \eA{Pao-yü put out his arm and pulled Hsi Jen. "You've been standing for so
        long," he smiled, "that you must be feeling tired."}
}
{
}

\Verse{135}
{
    \zA{一面说，一面拉他身旁 坐下。}
}
{
    \eA{Saying this, he dragged her down and made her take a seat next to him. "Here
        you've again forgotten!"}
}
{
}

\Verse{136}
{
    \zA{袭人笑道：“可是又忘了：趁宝姑娘在院子里，你和他说，烦他们莺儿来打 上几根绦子。”}
}
{
    \eA{laughingly exclaimed Hsi Jen. "Avail yourself now that Miss Pao-ch'ai is in
        the court to tell her to kindly bid their Ying Erh come and plait a few
        girdles with twisted cords." "How lucky it is you've reminded me?" Pao-yü
        observed with a smile.}
}
{
}

\Verse{137}
{
    \zA{宝玉笑道：“亏了你提起来。”}
}
{
    \eA{And putting, while he spoke, his head out of the window: "Cousin Pao-ch'ai,"
        he cried,}
}
{
}

\Verse{138}
{
    \zA{说着，便仰头向窗外道：“宝姐姐， 吃过饭叫莺儿来，烦他打几根绦子，可得闲儿？”}
}
{
    \eA{"when you've had your repast, do tell Ying Erh to come over. I would like to
        ask her to plait a few girdles for me. Has she got the time to spare?"
        Pao-ch'ai heard him speak; and turning round: "How about no time?"}
}
{
}

\Verse{139}
{
    \zA{宝钗听见，回头道：“是了，一 会儿就叫他来。”}
}
{
    \eA{she answered. "I'll tell her by and bye to come; it will be all right."
        Dowager lady Chia and the others, however, failed to catch distinctly the
        drift of their talk;}
}
{
}

\Verse{140}
{
    \zA{贾母等尚未听真，都止步问宝钗何事。}
}
{
    \eA{and they halted and made inquiries of Pao-ch'ai what it was about. Pao-ch'ai
        gave them the necessary explanations. "My dear child," remarked old lady
        Chia,}
}
{
}

\Verse{141}
{
    \zA{宝钗说明了，贾母便说道： “好孩子，你叫他来替你兄弟打几根罢。}
}
{
    \eA{"do let her come and twist a few girdles for your cousin! And should you be
        in need of any one for anything, I have over at my place a whole number of
        servant-girls doing nothing!}
}
{
}

\Verse{142}
{
    \zA{你要人使，我那里闲的丫头多着的呢。}
}
{
    \eA{Out of them, you are at liberty to send for any you like to wait on you!"
        "We'll send her to plait them!"}
}
{
}

\Verse{143}
{
    \zA{你 喜欢谁，只管叫来使唤。”}
}
{
    \eA{Mrs. Hsüeh and Pao-ch'ai observed smilingly with one consent. "What can we
        want her for?}
}
{
}

\Verse{144}
{
    \zA{薛姨妈宝钗等都笑道：“只管叫他来做就是了。}
}
{
    \eA{she also daily idles her time way and is up to every mischief!" But chatting
        the while, they were about to proceed on their way when they unexpectedly
        caught sight of Hsiang-yün,}
}
{
}

\Verse{145}
{
    \zA{有什么 使唤的去处!他天天也是闲着淘气。”}
}
{
    \eA{P'ing Erh, Hsiang Lin and other girls picking balsam flowers near the rocks;
        who, as soon as they saw the company approaching,}
}
{
}

\Verse{146}
{
    \zA{大家说着，往前正走，忽见湘云、平儿、香 菱等在山石边掐凤仙花呢，见了他们走来，都迎上来了。}
}
{
    \eA{advanced to welcome them. Shortly, they all sallied out of the garden.
        Madame Wang was worrying lest dowager lady Chia's strength might be
        exhausted, and she did her utmost to induce her to enter the drawing room
        and sit down.}
}
{
}

\Verse{147}
{
    \zA{少顷出至园外，王夫人恐贾母乏了，便欲让至上房内坐，贾母也觉脚酸，便点 头依允。}
}
{
    \eA{Old lady Chia herself was feeling her legs quite tired out, so she at once
        nodded her head and expressed her assent. Madame Wang then directed a
        waiting-maid to hurriedly precede them, and get ready the seats.}
}
{
}

\Verse{148}
{
    \zA{王夫人便命丫头忙先去铺设坐位。}
}
{
    \eA{But as Mrs. Chao had, about this time, pleaded indisposition, there was only
        therefore Mrs. Chou,}
}
{
}

\Verse{149}
{
    \zA{那时赵姨娘推病，只有周姨娘与那老婆 丫头们忙着打帘子，立靠背，铺褥子。}
}
{
    \eA{with the matrons and servant-girls at hand, so they had ample to do to raise
        the portières, to put the back-cushions in their places, and to spread out
        the rugs.}
}
{
}

\Verse{150}
{
    \zA{贾母扶着凤姐儿进来，与薛姨妈分宾主坐了， 宝钗湘云坐在下面。}
}
{
    \eA{Dowager lady Chia stepped into the room, leaning on lady Feng's arm. She and
        Mrs. Hsüeh took their places, with due regard to the distinction between
        hostess and visitors;}
}
{
}

\Verse{151}
{
    \zA{王夫人亲自捧了茶来，奉与贾母；李宫裁捧与薛姨妈。}
}
{
    \eA{and Hsüeh Pao-ch'ai and Shih Hsiang-yün seated themselves below. Madame Wang
        then came forward, and presented with her own hands tea to old lady Chia,}
}
{
}

\Verse{152}
{
    \zA{贾母向 王夫人道：“让他们小妯娌们伏侍罢，你在那里坐下，好说话儿。”}
}
{
    \eA{while Li Kung-ts'ai handed a cup to Mrs. Hsüeh. "You'd better let those
        young sisters-in law do the honours," remonstrated old lady Chia, "and sit
        over there so that we may be able to have a chat."}
}
{
}

\Verse{153}
{
    \zA{王夫人方向一 张小杌子上坐下，便吩咐凤姐儿道：“老太太的饭放在这里，添了东西来。”}
}
{
    \eA{Madame Wang at length sat on a small bench. "Let our worthy senior's
        viands," she cried, addressing herself to lady Feng, "be served here. And
        let a few more things be brought!" Lady Feng acquiesced without delay,}
}
{
}

\Verse{154}
{
    \zA{凤姐 儿答应出去，便命人去贾母那边告诉。}
}
{
    \eA{and she told a servant to cross over to their old mistress' quarters and to
        bid the matrons, employed in that part of the household,}
}
{
}

\Verse{155}
{
    \zA{那边的老婆们忙往外传了，丫头们忙都赶过 来。}
}
{
    \eA{promptly go out and summon the waiting-girls. The various waiting-maids
        arrived with all despatch. Madame Wang directed them to ask their young
        ladies round.}
}
{
}

\Verse{156}
{
    \zA{王夫人便命：“请姑娘们去。”}
}
{
    \eA{But after a protracted absence on the errand, only two of the girls turned
        up:}
}
{
}

\Verse{157}
{
    \zA{请了半天，只有探春惜春两个来了；迎春身上 不耐烦，不吃饭；那黛玉是不消说，十顿饭只好吃五顿，众人也不着意了。}
}
{
    \eA{T'an Ch'un and Hsi Ch'un. Ying Ch'un, was not, in her state of health, equal
        to the fatigue, or able to put anything in her mouth, and Lin Tai-yü,
        superfluous to add, could only safely partake of five out of ten meals,}
}
{
}

\Verse{158}
{
    \zA{少顷饭至，众人调放了桌子。}
}
{
    \eA{so no one thought anything of their non-appearance. Presently the eatables
        were brought,}
}
{
}

\Verse{159}
{
    \zA{凤姐儿用手巾裹了一把牙箸，站在地下，笑道： “老祖宗和姨妈不用让，还听我说就是了。”}
}
{
    \eA{and the servants arranged them in their proper places on the table. Lady
        Feng took a napkin and wrapped a bundle of chopsticks in it. "Venerable
        ancestor and you, Mrs. Hsüeh," she smiled, standing the while below,}
}
{
}

\Verse{160}
{
    \zA{贾母笑向薛姨妈道：“我们就是这样。”}
}
{
    \eA{"there's no need of any yielding! Just you listen to me and I'll make things
        all right." "Let's do as she wills!"}
}
{
}

\Verse{161}
{
    \zA{薛姨妈笑着应了。}
}
{
    \eA{old lady Chia remarked to Mrs. Hsüeh laughingly.}
}
{
}

\Verse{162}
{
    \zA{于是凤姐放下四双箸：上面两双是贾母薛姨妈，两边是宝钗湘云 的。}
}
{
    \eA{Mrs. Hsüeh signified her approval with a smile; so lady Feng placed, in due
        course, four pairs of chopsticks on the table; the two pairs on the upper
        end for dowager lady Chia and Mrs. Hsüeh;}
}
{
}

\Verse{163}
{
    \zA{王夫人李宫裁等都站在地下，看着放菜。}
}
{
    \eA{those on the two sides for Hsüeh Pao-ch'ai and Shih Hsiang-yün. Madame Wang,
        Li Kung-ts'ai and a few others,}
}
{
}

\Verse{164}
{
    \zA{凤姐先忙着要干净家伙来，替宝玉拣 菜。}
}
{
    \eA{stood together below and watched the attendants serve the viands. Lady Feng
        first and foremost hastily asked for clean utensils,}
}
{
}

\Verse{165}
{
    \zA{少顷，莲叶汤来了，贾母看过了，王夫人回头见玉钏儿在那里，便命玉钏儿与 宝玉送去。}
}
{
    \eA{and drew near the table to select some eatables for Pao-yü. Presently, the
        soup \_à la\_ lotus leaves arrived. After old lady Chia had well scrutinised
        it, Madame Wang turned her head, and catching sight of Yü Ch'uan-erh,}
}
{
}

\Verse{166}
{
    \zA{凤姐道：“他一个人难拿。”}
}
{
    \eA{she immediately commissioned her to take some over to Pao-yü.}
}
{
}

\Verse{167}
{
    \zA{可巧莺儿和同喜都来了，宝钗知道他们已 吃了饭，便向莺儿道：“宝二爷正叫你去打绦子，你们两个同去罢。”}
}
{
    \eA{"She can't carry it single-handed," demurred lady Feng. But by a strange
        coincidence, Ying Erh then walked into the room along with Hsi Erh, and
        Pao-ch'ai knowing very well that they had already had their meal forthwith
        said to Ying Erh: "Your Master Secundus, Mr. Pao-yü,}
}
{
}

\Verse{168}
{
    \zA{莺儿答应着， 和玉钏儿出来。}
}
{
    \eA{just asked that you should go and twist a few girdles for him; so you two
        might as well proceed together!"}
}
{
}

\Verse{169}
{
    \zA{莺儿道：“这么远，怪热的，那可怎么端呢？”}
}
{
    \eA{Ying Erh expressed her readiness and left the apartment, in company with Yü
        Ch'uan-erh. "How can you carry it, so very hot as it is,}
}
{
}

\Verse{170}
{
    \zA{玉钏儿笑道：“你 放心，我自有道理。”}
}
{
    \eA{the whole way there?" observed Ying Erh. "Don't distress yourself!" rejoined
        Yü Ch'uan smiling.}
}
{
}

\Verse{171}
{
    \zA{说着，便命一个婆子来，将汤饭等类放在一个捧盒里，命他 端了跟着，他两个却空着手走。}
}
{
    \eA{"I know how to do it." Saying this, she directed a matron to come and place
        the soup, rice and the rest of the eatables in a present box; and bidding
        her lay hold of it and follow them, the two girls sped on their way with
        empty hands,}
}
{
}

\Verse{172}
{
    \zA{一直到了怡红院门口，玉钏儿方接过来了，同着莺 儿进入房中。}
}
{
    \eA{and made straight for the entrance of the I Hung court. Here Yü Ch'uan-erh
        at length took the things herself, and entered the room in company with Ying
        Erh.}
}
{
}

\Verse{173}
{
    \zA{袭人、麝月、秋纹三个人正和宝玉玩笑呢，见他两个来了，都忙起来笑道：“你 们两个来的？}
}
{
    \eA{The trio, Hsi Jen, She Yüeh and Ch'iu Wen were at the time chatting and
        laughing with Pao-yü; but the moment they saw their two friends arrive they
        speedily jumped to their feet. "How is it," they exclaimed laughingly,}
}
{
}

\Verse{174}
{
    \zA{怎么碰巧一齐来了。”}
}
{
    \eA{"that you two drop in just the nick of time? Have you come together?"}
}
{
}

\Verse{175}
{
    \zA{一面说，一面接过来。}
}
{
    \eA{With these words on their lips, they descended to greet them.}
}
{
}

\Verse{176}
{
    \zA{玉钏儿便向一张杌子上 坐下；莺儿不敢坐，袭人便忙端了个脚踏来，莺儿还不敢坐。}
}
{
    \eA{Yü Ch'uan took at once a seat on a small stool. Ying Erh, however, did not
        presume to seat herself; and though Hsi Jen was quick enough in moving a
        foot-stool for her, Ying Erh did not still venture to sit down.}
}
{
}

\Verse{177}
{
    \zA{宝玉见莺儿来了，却 倒十分欢喜；见了玉钏儿，便想起他姐姐金钏儿来了，又是伤心，又是惭愧，便把 莺儿丢下，且和玉钏儿说话。}
}
{
    \eA{Ying Erh's arrival filled Pao-yü with intense delight. But as soon as he
        noticed Yü Ch'uan-erh, he recalled to memory her sister Chin Ch'uan-erh, and
        he felt wounded to the very heart, and overpowered with shame.}
}
{
}

\Verse{178}
{
    \zA{袭人见把莺儿不理，恐莺儿没好意思的，又见莺儿不 肯坐，便拉了莺儿出来，到那边屋里去吃茶说话儿去了。}
}
{
    \eA{And, without troubling his mind about Ying Erh, he addressed his remarks to
        Yü Ch'uan-erh. Hsi Jen saw very well that Ying Erh failed to attract his
        attention and she began to fear lest she felt uncomfortable; and when she
        further realised that Ying Erh herself would not take a seat,}
}
{
}

\Verse{179}
{
    \zA{这里麝月等预备了碗箸来伺候吃饭。}
}
{
    \eA{she drew her out of the room and repaired with her into the outer apartment,
        where they had a chat over their tea.}
}
{
}

\Verse{180}
{
    \zA{宝玉只是不吃，问玉钏儿道：“你母亲身 上好？”}
}
{
    \eA{She Yüeh and her companions had, in the meantime, got the bowls and
        chopsticks ready and came to wait upon (Pao-yü) during his meal.}
}
{
}

\Verse{181}
{
    \zA{玉钏儿满脸娇嗔，正眼也不看宝玉，半日方说了一个“好”字。}
}
{
    \eA{But Pao-yü would not have anything to eat. "Is your mother all right," he
        forthwith inquired of Yü Ch'uan-erh. An angry scowl crept over Yü
        Ch'uan-erh's face.}
}
{
}

\Verse{182}
{
    \zA{宝玉便觉 没趣，半日，只得又陪笑问道：“谁叫你替我送来的？”}
}
{
    \eA{She did not even look straight at Pao-yü. And only after a long pause was it
        that she at last uttered merely the words, "all right," by way of reply.}
}
{
}

\Verse{183}
{
    \zA{玉钏儿道：“不过是奶奶 太太们！”}
}
{
    \eA{Pao-yü, therefore, found talking to her of little zest. But after a
        protracted silence he felt impelled to again force a smile,}
}
{
}

\Verse{184}
{
    \zA{宝玉见他还是哭丧着脸，便知他是为金钏儿的原故。}
}
{
    \eA{and to ask: "Who told you to bring these things over to me?" "The ladies,"
        answered Yü Chuan-erh. Pao-yü discerned the mournful expression,}
}
{
}

\Verse{185}
{
    \zA{待要虚心下气哄他， 又见人多，不好下气的，因而便寻方法将人都支出去，然后又陪笑问长问短。}
}
{
    \eA{which still beclouded her countenance and he readily jumped at the
        conclusion that it must be entirely occasioned by the fate which had
        befallen Chin Ch'uan-erh, but when fain to put on a meek and unassuming
        manner, and endeavour to cheer her,}
}
{
}

\Verse{186}
{
    \zA{那玉 钏儿先虽不欲理他，只管见宝玉一些性气也没有，凭他怎么丧谤，还是温存和气， 自己倒不好意思的了，脸上方有三分喜色。}
}
{
    \eA{he saw how little he could demean himself in the presence of so many people,
        and consequently he did his best and discovered the means of getting every
        one out of the way. Afterwards, straining another smile, he plied her with
        all sorts of questions. Yü Ch'uan-erh, it is true, did not at first choose
        to heed his advances,}
}
{
}

\Verse{187}
{
    \zA{宝玉便笑央道：“好姐姐，你把那汤端 了来，我尝尝。”}
}
{
    \eA{yet when she observed that Pao-yü did not put on any airs, and, that in
        spite of all her querulous reproaches, he still continued pleasant and
        agreeable,}
}
{
}

\Verse{188}
{
    \zA{玉钏儿道：“我从不会喂人东西，等他们来了再喝。”}
}
{
    \eA{she felt disconcerted and her features at last assumed a certain expression
        of cheerfulness. Pao-yü thereupon smiled. "My dear girl," he said,}
}
{
}

\Verse{189}
{
    \zA{宝玉笑道： “我不是要你喂我，我因为走不动，你递给我喝了，你好赶早回去交代了，好吃饭 去。}
}
{
    \eA{as he gave way to entreaties, "bring that soup and let me taste it!" "I've
        never been in the habit of feeding people," Yü Ch'uan-erh replied. "You'd
        better wait till the others return; you can have some then."}
}
{
}

\Verse{190}
{
    \zA{我只管耽误了时候，岂不饿坏了你。}
}
{
    \eA{"I don't want you to feed me," laughed Pao-yü. "It's because I can't move
        about that I appeal to you.}
}
{
}

\Verse{191}
{
    \zA{你要懒怠动，我少不得忍着疼下去取去。”}
}
{
    \eA{Do let me have it! You'll then get back early and be able, when you've
        handed over the things,}
}
{
}

\Verse{192}
{
    \zA{说着，便要下床，扎挣起来，禁不住“嗳哟”之声。}
}
{
    \eA{to have your meal. But were I to go on wasting your time, won't you feel
        upset from hunger? Should you be lazy to budge,}
}
{
}

\Verse{193}
{
    \zA{玉钏儿见他这般，也忍不过， 起身说道：“躺下去罢!那世里造的孽，这会子现世现报，叫我那一个眼睛瞧的上！”}
}
{
    \eA{well then, I'll endure the pain and get down and fetch it myself." As he
        spoke, he tried to alight from bed. He strained every nerve, and raised
        himself, but unable to stand the exertion,}
}
{
}

\Verse{194}
{
    \zA{一面说，一面哧的一声又笑了，端过汤来。}
}
{
    \eA{he burst out into groans. At the sight of his anguish, Yü Ch'uan-erh had not
        the heart to refuse her help.}
}
{
}

\Verse{195}
{
    \zA{宝玉笑道：“好姐姐你要生气，只管在 这里生罢，见了老太太、太太，可和气着些。}
}
{
    \eA{Springing up, "Lie down!" she cried. "In what former existence did you
        commit such evil that your retribution in the present one is so apparent?
        Which of my eyes however can brook looking at you going on in that way?"}
}
{
}

\Verse{196}
{
    \zA{若还这样，你就要挨骂了。”}
}
{
    \eA{While taunting him, she again blurted out laughing, and brought the soup
        over to him.}
}
{
}

\Verse{197}
{
    \zA{玉钏儿 道：“吃罢，吃罢!你不用和我甜嘴蜜舌的了，我都知道啊！”}
}
{
    \eA{"My dear girl;" smiled Pao-yü, "if you want to show temper, better do so
        here! When you see our venerable senior and madame, my mother, you should be
        a little more even-tempered,}
}
{
}

\Verse{198}
{
    \zA{说着，催宝玉喝了 两口汤。}
}
{
    \eA{for if you still behave like this, you'll at once get a scolding!"}
}
{
}

\Verse{199}
{
    \zA{宝玉故意说不好吃。}
}
{
    \eA{"Eat away, eat away!" urged Yü Ch'uan-erh.}
}
{
}

\Verse{200}
{
    \zA{玉钏儿撇嘴道：“阿弥陀佛!这个还不好吃，也不知 什么好吃呢！”}
}
{
    \eA{"There's no need for you to be so sweet-mouthed and honey-tongued with me. I
        don't put any faith in such talk!" So speaking, she pressed Pao-yü until he
        had two mouthfuls of soup.}
}
{
}

\Verse{201}
{
    \zA{宝玉道：“一点味儿也没有，你不信尝一尝，就知道了。”}
}
{
    \eA{"It isn't nice, it isn't nice!" Pao-yü purposely exclaimed. "Omi-to-fu!"
        ejaculated Yü Ch'uan-erh. "If this isn't nice, what's nice?"}
}
{
}

\Verse{202}
{
    \zA{玉钏儿 果真赌气尝了一尝。}
}
{
    \eA{"There's no flavour about it at all," resumed Pao-yü. "If you don't believe
        me taste it,}
}
{
}

\Verse{203}
{
    \zA{宝玉笑道：“这可好吃了！”}
}
{
    \eA{and you'll find out for yourself." Yü Ch'uan-erh in a tantrum actually put
        some of it to her lips.}
}
{
}

\Verse{204}
{
    \zA{玉钏儿听说，方解过他的意思来， 原是宝玉哄他喝一口，便说道：“你既说不喝，这会子说好吃，也不给你喝了。”}
}
{
    \eA{"Well," laughed Pao-yü, "it is nice!" This exclamation eventually enabled Yü
        Ch'uan to see what Pao-yü was driving at, for Pao-yü had in fact been trying
        to beguile her to have a mouthful. "As, at one moment, you say you don't
        want any," she forthwith observed, "and now you say it is nice,}
}
{
}

\Verse{205}
{
    \zA{宝玉只管陪笑央求要喝，玉钏儿又不给他，一面又叫人打发吃饭。}
}
{
    \eA{I won't give you any." While Pao-yü returned her smiles, he kept on
        earnestly entreating her to let him have some. Yü Ch'uan-erh however would
        still not give him any;}
}
{
}

\Verse{206}
{
    \zA{丫头方进来时，忽有人来回话，说：“傅二爷家的两个嬷嬷来请安，来见二爷。”}
}
{
    \eA{and she, at the same time, called to the servants to fetch what there was
        for him to eat. But the instant the waiting-maid put her foot into the room,
        servants came quite unexpectedly to deliver a message.}
}
{
}

\Verse{207}
{
    \zA{宝玉听说，便知是通判傅试家的嬷嬷来了。}
}
{
    \eA{"Two nurses," they said, "have arrived from the household of Mr. Fu,
        Secundus, to present his compliments.}
}
{
}

\Verse{208}
{
    \zA{那傅试原是贾政的门生，原来都赖贾家 的名声得意，贾政也着实看待，与别的门生不同；他那里常遣人来走动。}
}
{
    \eA{They have now come to see you, Mr. Secundus." As soon as Pao-yü heard this
        report, he felt sure that they must be nurses sent over from the household
        of Deputy Sub-Prefect, Fu Shih. This Fu Shih had originally been a pupil of
        Chia Cheng,}
}
{
}

\Verse{209}
{
    \zA{宝玉素昔 最厌勇男蠢妇的，今日却如何又命这两个婆子进来？}
}
{
    \eA{and had, indeed, had to rely entirely upon the reputation enjoyed by the
        Chia family for the realisation of his wishes. Chia Cheng had, likewise,}
}
{
}

\Verse{210}
{
    \zA{其中原来有个原故。}
}
{
    \eA{treated him with such genuine regard, and so unlike any of his other pupils,}
}
{
}

\Verse{211}
{
    \zA{只因那宝 玉闻得傅试有个妹子，名唤傅秋芳，也是个琼闺秀玉，常听人说才貌俱全，虽自未 亲睹，然遐思遥爱之心十分诚敬。}
}
{
    \eA{that he (Fu Shih) ever and anon despatched inmates from his mansion to come
        and see him so as to keep up friendly relations. Pao-yü had at all times
        entertained an aversion for bold-faced men and unsophisticated women, so why
        did he once more, on this occasion, issue directions that the two matrons
        should be introduced into his presence?}
}
{
}

\Verse{212}
{
    \zA{不命他们进来，恐薄了傅秋芳，因此连忙命让进 来。}
}
{
    \eA{There was, in fact, a reason for his action. It was simply that Pao-yü had
        come to learn that Fu Shih had a sister, Ch'iu-fang by name, a girl as
        comely as a magnificent gem,}
}
{
}

\Verse{213}
{
    \zA{那傅试原是暴发的，因傅秋芳有几分姿色，聪明过人，那傅试安心仗着妹子， 要与豪门贵族结亲，不肯轻意许人，所以耽误到如今。}
}
{
    \eA{and perfection itself, the report of outside people went, as much in
        intellect as in beauty. He had, it is true, not yet seen anything of her
        with his own eyes, but the sentiments, which made him think of her and
        cherish her, from a distance, were characterised by such extreme sincerity,}
}
{
}

\Verse{214}
{
    \zA{目今傅秋芳已二十三岁，尚 未许人。}
}
{
    \eA{that dreading lest he should, by refusing to admit the matrons, reflect
        discredit upon Fu Ch'iu-fang,}
}
{
}

\Verse{215}
{
    \zA{怎奈那些豪门贵族又嫌他本是穷酸，根基浅薄，不肯求配。}
}
{
    \eA{he was prompted to lose no time in expressing a wish that they should be
        ushered in. This Fu Shih had really risen from the vulgar herd,}
}
{
}

\Verse{216}
{
    \zA{那傅试与贾家 亲密，也自有一段心事。}
}
{
    \eA{so seeing that Ch'iu-fang possessed several traits of beauty and exceptional
        intellectual talents,}
}
{
}

\Verse{217}
{
    \zA{今日遣来的两个婆子，偏偏是极无知识的，闻得宝玉要见，进来只刚问了好， 说了没两句话。}
}
{
    \eA{Fu Shih arrived at the resolution of making his sister the means of joining
        relationship with the influential family of some honourable clan. And so
        unwilling was he to promise her lightly to any suitor that things were
        delayed up to this time. Therefore Fu Ch'iu-fang, though at present past her
        twentieth birthday,}
}
{
}

\Verse{218}
{
    \zA{那玉钏儿见生人来，也不和宝玉厮闹了，手里端着汤，却只顾听。}
}
{
    \eA{was not as yet engaged. But the various well-to-do families, belonging to
        honourable clans, looked down, on the other hand, on her poor and mean
        extraction,}
}
{
}

\Verse{219}
{
    \zA{宝玉又只顾和婆子说话，一面吃饭，伸手去要汤，两个人的眼睛都看着人，不想伸 猛了手，便将碗撞翻，将汤泼了宝玉手上。}
}
{
    \eA{holding her in such light esteem, as not to relish the idea of making any
        offer for her hand. So if Fu Shih cultivated intimate terms with the Chia
        household, he, needless to add, did so with an interested motive. The two
        matrons, deputed on the present errand, completely lacked, as it happened,}
}
{
}

\Verse{220}
{
    \zA{玉钏儿倒不曾烫着，吓了一跳，忙笑着： “这是怎么了？”}
}
{
    \eA{all knowledge of the world, and the moment they heard that Pao-yü wished to
        see them, they wended their steps inside. But no sooner had they inquired
        how he was,}
}
{
}

\Verse{221}
{
    \zA{慌的丫头们忙上来接碗。}
}
{
    \eA{and passed a few remarks than Yü Ch'uan-erh, becoming conscious of the
        arrival of strangers,}
}
{
}

\Verse{222}
{
    \zA{宝玉自己烫了手，倒不觉的，只管问玉 钏儿：“烫了那里了？}
}
{
    \eA{did not bandy words with Pao-yü, but stood with the plate of soup in her
        hands, engrossed in listening to the conversation. Pao-yü, again, was
        absorbed in speaking to the matrons;}
}
{
}

\Verse{223}
{
    \zA{疼不疼？”}
}
{
    \eA{and, while eating some rice,}
}
{
}

\Verse{224}
{
    \zA{玉钏儿和众人都笑了。}
}
{
    \eA{he stretched out his arm to get at the soup; but both his and her (Yü
        Ch'uan-erh's) eyes were rivetted on the women,}
}
{
}

\Verse{225}
{
    \zA{玉钏儿道：“你自己烫了， 只管问我。”}
}
{
    \eA{and as he thoughtlessly jerked out his hand with some violence, he struck
        the bowl and turned it clean over.}
}
{
}

\Verse{226}
{
    \zA{宝玉听了，方觉自己烫了。}
}
{
    \eA{The soup fell over Pao-yü's hand. But it did not hurt Yü Ch'uan-erh.}
}
{
}

\Verse{227}
{
    \zA{众人上来，连忙收拾。}
}
{
    \eA{She sustained, however, such a fright that she gave a start.}
}
{
}

\Verse{228}
{
    \zA{宝玉也不吃饭了， 洗手吃茶，又和那两个婆子说了两句话，然后两个婆子告辞出去。}
}
{
    \eA{"How did this happen!" she smilingly shouted with vehemence to the intense
        consternation of the waiting-maids, who rushed up and clasped the bowl. But
        notwithstanding that Pao-yü had scalded his own hand,}
}
{
}

\Verse{229}
{
    \zA{晴雯等送至桥边 方回。}
}
{
    \eA{he was quite unconscious of the accident;}
}
{
}

\Verse{230}
{
    \zA{那两个婆子见没人了，一行走一行谈论。}
}
{
    \eA{so much so, that he assailed Yü Ch'uan-erh with a heap of questions, as to
        where she had been burnt, and whether it was sore or not.}
}
{
}

\Verse{231}
{
    \zA{这一个笑道：“怪道有人说他们家的 宝玉是相貌好里头糊涂，中看不中吃，果然竟有些呆气。}
}
{
    \eA{Yü Ch'uan-erh and every one present were highly amused. "You yourself,"
        observed Yü Ch'uan-erh, "have been scalded, and do you keep on asking about
        myself?" At these words, Pao-yü became at last aware of the injury he had
        received.}
}
{
}

\Verse{232}
{
    \zA{他自己烫了手，倒问别人 疼不疼，这可不是呆了吗！”}
}
{
    \eA{The servants rushed with all promptitude and cleared the mess. But Pao-yü
        was not inclined to touch any more food. He washed his hands, drank a cup of
        tea,}
}
{
}

\Verse{233}
{
    \zA{那个又笑道：“我前一回来，还听见他家里许多人说， 千真万真有些呆气：大雨淋的水鸡儿似的，他反告诉别人：‘下雨了，快避雨去罢。’}
}
{
    \eA{and then exchanged a few further sentences with the two matrons. But
        subsequently, the two women said good-bye and quitted the room. Ch'ing Wen
        and some other girls saw them as far as the bridge, after which, they
        retraced their steps.}
}
{
}

\Verse{234}
{
    \zA{你说可笑不可笑？}
}
{
    \eA{The two matrons perceived, that there was no one about,}
}
{
}

\Verse{235}
{
    \zA{时常没人在跟前，就自哭自笑的，看见燕子就和燕子说话，河里 看见了鱼就和鱼儿说话，见了星星月亮，他不是长吁短叹的，就是咕咕哝哝的。}
}
{
    \eA{and while proceeding on their way, they started a conversation. "It isn't
        strange," smiled the one, "if people say that this Pao-yü of theirs is
        handsome in appearance,}
}
{
}

\Verse{236}
{
    \zA{且 一点刚性儿也没有，连那些毛丫头的气都受到了。}
}
{
    \eA{but stupid as far as brains go. Nice enough a thing to look at but not to
        put to one's lips; rather idiotic in fact; for he burns his own hand,}
}
{
}

\Verse{237}
{
    \zA{爱惜起东西来，连个线头儿都是 好的；遭塌起来，那怕值千值万都不管了。”}
}
{
    \eA{and then he asks some one else whether she's sore or not. Now, isn't this
        being a regular fool?" "The last time I came," the other remarked, also
        smiling, "I heard that many inmates of his family feel ill-will against him.}
}
{
}

\Verse{238}
{
    \zA{两个人一面说，一面走出园来回去， 不在话下。}
}
{
    \eA{In real truth he is a fool! For there he drips in the heavy downpour like a
        water fowl, and instead of running to shelter himself,}
}
{
}

\Verse{239}
{
    \zA{且说袭人见人去了，便携了莺儿过来问宝玉：“打什么绦子？”}
}
{
    \eA{he reminds other people of the rain, and urges them to get quick out of the
        wet. Now, tell me, isn't this ridiculous, eh? Time and again, when no one is
        present,}
}
{
}

\Verse{240}
{
    \zA{宝玉笑向莺儿 道：“才只顾说话，就忘了你了。}
}
{
    \eA{he cries to himself, then laughs to himself. When he sees a swallow, he
        instantly talks to it; when he espies a fish, in the river,}
}
{
}

\Verse{241}
{
    \zA{烦你来不为别的，替我打几根络子。”}
}
{
    \eA{he forthwith speaks to it. At the sight of stars or the moon, if he doesn't
        groan and sigh,}
}
{
}

\Verse{242}
{
    \zA{莺儿道： “装什么的络子？”}
}
{
    \eA{he mutters and mutters. Indeed, he hasn't the least bit of character;}
}
{
}

\Verse{243}
{
    \zA{宝玉见问，便笑道：“不管装什么的，你都每样打几个罢。”}
}
{
    \eA{so much so, that he even puts up with the temper shown by those low-bred
        maids. If he takes a fancy to a thing, it's nice enough even though it be a
        bit of thread.}
}
{
}

\Verse{244}
{
    \zA{莺儿拍手笑道：“这还了得，要这样，十年也打不完了。”}
}
{
    \eA{But as for waste, what does he mind? A thing may be worth a thousand or ten
        thousand pieces of money, he doesn't worry his mind in the least about it."}
}
{
}

\Verse{245}
{
    \zA{宝玉笑道：“好姑娘， 你闲着也没事，都替我打了罢。”}
}
{
    \eA{While they talked, they reached the exterior of the garden, and they betook
        themselves back to their home; where we will leave them.}
}
{
}

\Verse{246}
{
    \zA{袭人笑道：“那里一时都打的完？}
}
{
    \eA{As soon as Hsi Jen, for we will return to her, saw the women leave the room,
        she took Ying Erh by the hand and led her in,}
}
{
}

\Verse{247}
{
    \zA{如今先拣要紧 的打几个罢。”}
}
{
    \eA{and they asked Pao-yü what kind of girdle he wanted made. "I was just now so
        bent upon talking,"}
}
{
}

\Verse{248}
{
    \zA{莺儿道：“什么要紧，不过是扇子，香坠儿，汗巾子。”}
}
{
    \eA{Pao-yü smiled to Ying Erh, "that I forgot all about you. I put you to the
        trouble of coming, not for anything else, but that you should also make me a
        few nets."}
}
{
}

\Verse{249}
{
    \zA{宝玉道： “汗巾子就好。”}
}
{
    \eA{"Nets! To put what in?" Ying Erh inquired. Pao-yü, at this question,}
}
{
}

\Verse{250}
{
    \zA{莺儿道：“汗巾子是什么颜色？”}
}
{
    \eA{put on a smile. "Don't concern yourself about what they are for!" he
        replied. "Just make me a few of each kind!"}
}
{
}

\Verse{251}
{
    \zA{宝玉道：“大红的。”}
}
{
    \eA{Ying Erh clapped her hand and laughed. "Could this ever be done!"}
}
{
}

\Verse{252}
{
    \zA{莺儿道： “大红的须是黑络子才好看，或是石青的，才压得住颜色。”}
}
{
    \eA{she cried, "If you want all that lot, why, they couldn't be finished in ten
        years time." "My dear girl," smiled Pao-yü, "work at them for me then
        whenever you are at leisure,}
}
{
}

\Verse{253}
{
    \zA{宝玉道：“松花色配 什么？”}
}
{
    \eA{and have nothing better to do." "How could you get through them all in a
        little time?"}
}
{
}

\Verse{254}
{
    \zA{莺儿道：“松花配桃红。”}
}
{
    \eA{Hsi Jen interposed smilingly. "First choose now therefore such as are most
        urgently needed and make a couple of them."}
}
{
}

\Verse{255}
{
    \zA{宝玉笑道：“这才娇艳。}
}
{
    \eA{"What about urgently needed?" Ying-Erh exclaimed, "They are merely used for
        fans,}
}
{
}

\Verse{256}
{
    \zA{再要雅淡之中带些娇 艳。”}
}
{
    \eA{scented pendants and handkerchiefs." "Nets for handkerchiefs will do all
        right."}
}
{
}

\Verse{257}
{
    \zA{莺儿道：“葱绿柳黄可倒还雅致。”}
}
{
    \eA{Pao-yü answered. "What's the colour of your handkerchief?"}
}
{
}

\Verse{258}
{
    \zA{宝玉道：“也罢了。}
}
{
    \eA{inquired Ying Erh. "It's a deep red one."}
}
{
}

\Verse{259}
{
    \zA{也打一条桃红，再 打一条葱绿。”}
}
{
    \eA{Pao-yü rejoined. "For a deep red one," continued Ying Erh, "a black net will
        do very nicely,}
}
{
}

\Verse{260}
{
    \zA{莺儿道：“什么花样呢？”}
}
{
    \eA{or one of dark green. Both these agree with the colour."}
}
{
}

\Verse{261}
{
    \zA{宝玉道：“也有几样花样？”}
}
{
    \eA{"What goes well with brown?" Pao-yü asked. "Peach-red goes well with brown."}
}
{
}

\Verse{262}
{
    \zA{莺儿道： “‘一炷香’，‘朝天凳’，‘象眼块’，‘方胜’，‘连环’，‘梅花’，‘柳 叶’。”}
}
{
    \eA{Ying Erh added. "That will make them look gaudy!" Pao-yü observed. "Yet with
        all their plainness, they should be somewhat gaudy." "Leek-green and
        willow-yellow are what are most to my taste,"}
}
{
}

\Verse{263}
{
    \zA{宝玉道：“前儿你替三姑娘打的那花样是什么？”}
}
{
    \eA{Ying Erh pursued. "Yes, they'll also do!" Pao-yü retorted. "But make one of
        peach-red too and then one of leek-green."}
}
{
}

\Verse{264}
{
    \zA{莺儿道：“是‘攒心梅 花’。”}
}
{
    \eA{"Of what design?" Ying Erh remarked. "How many kinds of designs are there?"}
}
{
}

\Verse{265}
{
    \zA{宝玉道：“就是那样好。”}
}
{
    \eA{Pao-yü said. "There are 'the stick of incense,' 'stools upset towards
        heaven,'}
}
{
}

\Verse{266}
{
    \zA{一面说，一面袭人刚拿了线来。}
}
{
    \eA{'part of elephant's eyes,' 'squares,' 'chains,' 'plum blossom,'}
}
{
}

\Verse{267}
{
    \zA{窗外婆子说： “姑娘们的饭都有了。”}
}
{
    \eA{and 'willow leaves." Ying Erh answered. "What was the kind of design you
        made for Miss Tertia the other day?"}
}
{
}

\Verse{268}
{
    \zA{宝玉道：“你们吃饭去，快吃了来罢。”}
}
{
    \eA{Pao-yü inquired. "It was the 'plum blossom with piled cores,'" Ying Erh
        explained in reply.}
}
{
}

\Verse{269}
{
    \zA{袭人笑道：“有 客在这里。}
}
{
    \eA{"Yes, that's nice." Pao-yü rejoined. As he uttered this remark, Hsi Jen
        arrived with the cords.}
}
{
}

\Verse{270}
{
    \zA{我们怎么好意思去呢？”}
}
{
    \eA{But no sooner were they brought than a matron cried, from outside the
        window:}
}
{
}

\Verse{271}
{
    \zA{莺儿一面理线，一面笑道：“这打那里说起？}
}
{
    \eA{"Girls, your viands are ready!" "Go and have your meal," urged Pao-yü, "and
        come back quick after you've had it."}
}
{
}

\Verse{272}
{
    \zA{正经快吃去罢。”}
}
{
    \eA{"There are visitors here,"}
}
{
}

\Verse{273}
{
    \zA{袭人等听说，方去了，只留下两个小丫头呼唤。}
}
{
    \eA{Hsi Jen smiled, "and how can I very well go?" "What makes you say so?" Ying
        Erh laughed, while adjusting the cords.}
}
{
}

\Verse{274}
{
    \zA{宝玉一面看莺儿打络子，一面说闲话。}
}
{
    \eA{"It's only right and proper that you should go and have your food at once
        and then return."}
}
{
}

\Verse{275}
{
    \zA{因问他：“十几岁了？”}
}
{
    \eA{Hearing this, Hsi Jen and her companions went off,}
}
{
}

\Verse{276}
{
    \zA{莺儿手里打着， 一面答话：“十五岁了。”}
}
{
    \eA{leaving behind only two youthful servant-girls to answer the calls. Pao-yü
        watched Ying Erh make the nets.}
}
{
}

\Verse{277}
{
    \zA{宝玉道：“你本姓什么？”}
}
{
    \eA{But, while keeping his eyes intent on her,}
}
{
}

\Verse{278}
{
    \zA{莺儿道：“姓黄。”}
}
{
    \eA{he talked at the same time of one thing and then another,}
}
{
}

\Verse{279}
{
    \zA{宝玉笑 道：“这个姓名倒对了，果然是个‘黄莺儿’。”}
}
{
    \eA{and next went on to ask her how far she was in her teens. Ying Erh continued
        plaiting. "I'm sixteen," she simultaneously rejoined.}
}
{
}

\Verse{280}
{
    \zA{莺儿笑道：“我的名字本来是两 个字，叫做金莺，姑娘嫌拗口，只单叫莺儿，如今就叫开了。”}
}
{
    \eA{"What was your original surname?" Pao-yü added. "It was Huang;" answered
        Ying Erh. "That's just the thing," Pao-yü smiled; "for in real truth there's
        the 'Huang Ying-erh;' (oriole)." "My name, at one time,}
}
{
}

\Verse{281}
{
    \zA{宝玉道：“宝姐姐 也就算疼你了。}
}
{
    \eA{consisted of two characters," continued Ying Erh. "I was called Chin Ying;}
}
{
}

\Verse{282}
{
    \zA{明儿宝姐姐出嫁，少不得是你跟了去了。”}
}
{
    \eA{but Miss Pao-ch'ai didn't like it, as it was difficult to pronounce, and
        only called me Ying Erh;}
}
{
}

\Verse{283}
{
    \zA{莺儿抿嘴一笑。}
}
{
    \eA{so now I've come to be known under that name."}
}
{
}

\Verse{284}
{
    \zA{宝玉笑 道：“我常常和你花大姐姐说，明儿也不知那一个有造化的消受你们主儿两个呢。”}
}
{
    \eA{"One can very well say that cousin Pao-ch'ai is fond of you!" Pao-yü
        pursued. "By and bye, when she gets married, she's sure to take you along
        with her." Ying Erh puckered up her lips, and gave a significant smile.}
}
{
}

\Verse{285}
{
    \zA{莺儿笑道：“你还不知我们姑娘，有几样世上的人没有的好处呢，模样儿还在其次。”}
}
{
    \eA{"I've often told Hsi Jen," Pao-yü smiled, "that I can't help wondering
        who'll shortly be the lucky ones to win your mistress and yourself." "You
        aren't aware," laughed Ying Erh, "that our young mistress possesses several
        qualities not to be found in a single person in this world;}
}
{
}

\Verse{286}
{
    \zA{宝玉见莺儿娇腔婉转，语笑如痴，早不胜其情了，那堪更提起宝钗来？}
}
{
    \eA{her face is a second consideration." Pao-yü noticed how captivating Ying
        Erh's tone of voice was, how complaisant she was, and how simpleton-like
        unaffected in her language and smiles,}
}
{
}

\Verse{287}
{
    \zA{便问道：“什 么好处？}
}
{
    \eA{and he soon felt the warmest affection for her;}
}
{
}

\Verse{288}
{
    \zA{你细细儿的告诉我听。”}
}
{
    \eA{and particularly so, when she started the conversation about Pao-ch'ai.}
}
{
}

\Verse{289}
{
    \zA{莺儿道：“我告诉你，你可不许告诉他。”}
}
{
    \eA{"Where do her qualities lie?" he readily inquired. "My dear girl, please
        tell me!" "If I tell you,"}
}
{
}

\Verse{290}
{
    \zA{宝玉 笑道：“这个自然。”}
}
{
    \eA{said Ying Erh, "you must, on no account, let her know anything about it
        again."}
}
{
}

\Verse{291}
{
    \zA{正说着，只听见外头说道：“怎么这么静悄悄的？”}
}
{
    \eA{"This goes without saying," smiled Pao-yü. But this answer was still on his
        lips, when they overheard some one outside remark:}
}
{
}

\Verse{292}
{
    \zA{二人回头看时，不是别人， 正是宝钗来了。}
}
{
    \eA{"How is it that everything is so quiet?" Both gazed round to see who
        possibly it could be.}
}
{
}

\Verse{293}
{
    \zA{宝玉忙让坐。}
}
{
    \eA{They discovered, strange enough,}
}
{
}

\Verse{294}
{
    \zA{宝钗坐下，因问莺儿：“打什么呢？”}
}
{
    \eA{no one else than Pao-ch'ai herself. Pao-yü hastily offered her a seat.
        Pao-ch'ai seated herself,}
}
{
}

\Verse{295}
{
    \zA{一面问，一面 向他手里去瞧，才打了半截儿。}
}
{
    \eA{and then wanted to know what Ying Erh was busy plaiting. Inquiring the
        while, she approached her and scrutinised what she held in her hands,}
}
{
}

\Verse{296}
{
    \zA{宝钗笑道：“这有什么趣儿，倒不如打个络子把玉 络上呢。”}
}
{
    \eA{half of which had by this time been done. "What's the fun of a thing like
        this?" she said. "Wouldn't it be preferable to plait a net, and put the jade
        in it?"}
}
{
}

\Verse{297}
{
    \zA{一句话提醒了宝玉，便拍手笑道：“倒是姐姐说的是，我就忘了。}
}
{
    \eA{This allusion suggested the idea to Pao-yü. Speedily clapping his hands, he
        smiled and exclaimed: "Your idea is splendid, cousin.}
}
{
}

\Verse{298}
{
    \zA{只是 配个什么颜色才好？”}
}
{
    \eA{I'd forgotten all about it! The only thing is what colour will suit it
        best?"}
}
{
}

\Verse{299}
{
    \zA{宝钗道：“用鸦色断然使不得，大红又犯了色。}
}
{
    \eA{"It will never do to use mixed colours," Pao-ch'ai rejoined. "Deep red will,
        on one hand, clash with the colour;}
}
{
}

\Verse{300}
{
    \zA{黄的又不起 眼，黑的太暗。}
}
{
    \eA{while yellow is not pleasing to the eye; and black, on the other hand,}
}
{
}

\Verse{301}
{
    \zA{依我说，竟把你的金线拿来配着黑珠儿线，一根一根的拈上，打成 络子，那才好看。”}
}
{
    \eA{is too sombre. But wait, I'll try and devise something. Bring that gold cord
        and use it with the black beaded cord; and if you twist one of each
        together, and make a net with them,}
}
{
}

\Verse{302}
{
    \zA{宝玉听说，喜之不尽，一叠连声就叫袭人来取金线。}
}
{
    \eA{it will look very pretty!" Upon hearing this, Pao-yü was immeasurably
        delighted, and time after time he shouted to the servants to fetch the gold
        cord.}
}
{
}

\Verse{303}
{
    \zA{正值袭人端了两碗菜走进来，告诉宝玉道：“今儿奇怪，刚才太太打发人给我 送了两碗菜来。”}
}
{
    \eA{But just at that moment Hsi Jen stepped in, with two bowls of eatables. "How
        very strange this is to-day!" she said to Pao-yü. "Why, a few minutes back,
        my mistress, your mother, sent some one to bring me two bowls of viands."}
}
{
}

\Verse{304}
{
    \zA{宝玉笑道：“必定是今儿菜多，送给你们大家吃的。”}
}
{
    \eA{"The supply," replied Pao-yü smiling, "must have been so plentiful to-day,
        that they've sent some to every one of you."}
}
{
}

\Verse{305}
{
    \zA{袭人道： “不是，说指名给我的，还不叫过去磕头，这可是奇了。”}
}
{
    \eA{"It isn't that," continued Hsi Jen, "for they were distinctly given to me by
        name. What's more, I wasn't bidden go and knock my head; so this is indeed
        remarkable!"}
}
{
}

\Verse{306}
{
    \zA{宝钗笑道：“给你的你 就吃去，这有什么猜疑的。”}
}
{
    \eA{"If they're given to you," Pao-yü smiled, "why, you had better go and eat
        them. What's there in this to fill you with conjectures?"}
}
{
}

\Verse{307}
{
    \zA{袭人道：“从来没有的事，倒叫我不好意思的。”}
}
{
    \eA{"There's never been anything like this before," Hsi Jen added, "so, it makes
        me feel uneasy."}
}
{
}

\Verse{308}
{
    \zA{宝 钗抿嘴一笑，说道：“这就不好意思了？}
}
{
    \eA{Pao-ch'ai compressed her lips. "If this," she laughed; "makes you feel
        uneasy,}
}
{
}

\Verse{309}
{
    \zA{明儿还有比这个更叫你不好意思的呢！”}
}
{
    \eA{there will be by and bye other things to make you far more uneasy."}
}
{
}

\Verse{310}
{
    \zA{袭人听了话内有因，素知宝钗不是轻嘴薄舌奚落人的，自己想起上日王夫人的意思 来，便不再提了。}
}
{
    \eA{Hsi Jen realised that she implied something by her insinuations, as she knew
        from past experience that Pao-ch'ai was not one given to lightly and
        contemptuously poking fun at people; and, remembering the notions
        entertained by Madame Wang on the last occasion she had seen her, she
        dropped at once any further allusions to the subject and brought the
        eatables up to Pao-yü for his inspection.}
}
{
}

\Verse{311}
{
    \zA{将菜给宝玉看了，说：“洗了手来拿线。”}
}
{
    \eA{"I shall come and hold the cords," she observed, "as soon as I've rinsed my
        hands."}
}
{
}

\Verse{312}
{
    \zA{说毕，便一直出去了。}
}
{
    \eA{This said, she immediately quitted the apartment.}
}
{
}

\Verse{313}
{
    \zA{吃过饭洗了手进来，拿金线给莺儿打络子。}
}
{
    \eA{After her meal, she washed her hands and came inside to hold the gold cords
        for Ying Erh to plait the net with.}
}
{
}

\Verse{314}
{
    \zA{此时宝钗早被薛蟠遣人来请出去了。}
}
{
    \eA{By this time, Pao-ch'ai had been called away by a servant, despatched by
        Hsüeh P'an.}
}
{
}

\Verse{315}
{
    \zA{这里宝玉正看着打络子，忽见邢夫人那边遣了两个丫头送了两样果子来给他 吃，问他：“可走得了么？}
}
{
    \eA{But while Pao-yü was watching the net that was being made he caught sight,
        at a moment least expected, of two servant-girls, who came from the part of
        Madame Hsing of the other mansion, to bring him a few kinds of fruits, and
        to inquire whether he was able to walk.}
}
{
}

\Verse{316}
{
    \zA{要走的动，叫哥儿明儿过去散散心，太太着实惦记着呢。”}
}
{
    \eA{"If you can go about," they told him, "(our mistress) desires you, Mr.
        Pao-yü, to cross over to-morrow and have a little distraction. Her ladyship
        really longs to see you."}
}
{
}

\Verse{317}
{
    \zA{宝玉忙道：“要走得了，必定过来请太太的安去。}
}
{
    \eA{"Were I able to walk," Pao-yü answered with alacrity, "I would feel it my
        duty to go and pay my respects to your mistress!}
}
{
}

\Verse{318}
{
    \zA{疼的比先好些，请太太放心罢。”}
}
{
    \eA{Anyhow, the pain is better than before, so request your lady to allay her
        solicitude."}
}
{
}

\Verse{319}
{
    \zA{一面叫他两个坐下，一面又叫：“秋纹来，把才那果子拿一半送给林姑娘去。”}
}
{
    \eA{As he bade them both sit down, he, at the same time, called Ch'iu Wen.
        "Take," he said to her, "half of the fruits, just received, to Miss Lin as a
        present." Ch'iu Wen signified her obedience,}
}
{
}

\Verse{320}
{
    \zA{秋 纹答应了，刚欲去时，只听黛玉在院内说话。}
}
{
    \eA{and was about to start on her errand, when she heard Tai-yü talking in the
        court, and Pao-yü eagerly shout out:}
}
{
}

\Verse{321}
{
    \zA{宝玉忙叫快请。}
}
{
    \eA{"Request her to walk in at once!"}
}
{
}

\Verse{322}
{
    \zA{要知端底，且看下回分解。}
}
{
    \eA{But should there be any further particulars, which you, reader, might feel
        disposed to know,}
}
{
}

\Verse{323}
{
    \zA{(本章完)}
}
{
    \eA{peruse the details given in the following chapter.}
}
{
}
